% 第 8 章 PXL（前半：p175--210）
\bichapter{PXL}{PXL}
\label{chap:pxl}

本章介绍可编程可扩展语言 PXL，以及在 IC Validator 工具中函数定义与函数使用的基础。

PXL 包括以下各节：
\begin{itemize}
  \item PXL 体系结构（PXL Architecture）
  \item 一般定义（General Definitions）
  \item 变量与运算符（Variables and Operators）
  \item 数据类型与类型系统（Data Types and Typing）
  \item 流程控制（Flow Control）
  \item 函数（Functions）
  \item 程序结构（Program Structure）
  \item 模块化代码示例（Example of Modularized Code）
\end{itemize}

% ============================================================
\section{PXL 体系结构}
\subsection*{PXL Architecture}

PXL 完全可编程，具备下列一般可编程性特性：
\begin{itemize}
  \item 变量（Variables）
  \item 流程控制（Flow Control）
  \item 带标准索引的列表（Lists / arrays）
  \item 结构体（Structures）
  \item 哈希表（Hashes）
  \item 枚举类型（Enumerated types）
  \item 用户定义函数（User-defined functions）
  \item 数学、逻辑与字符串运算
  \item 支持命令行选项的宏（Macros）
  \item 环境变量（Environment variables）
  \item 静态作用域（Static scope）
\end{itemize}

这些特性可带来：
\begin{itemize}
  \item 直观的 runset 开发
  \item 易于维护的 runset
  \item 简洁的代码
  \item 通过子程序实现代码复用
  \item 支持用户定义函数
\end{itemize}

图~\ref{fig:pxl-structure} 给出 PXL 的基本结构。

\figplaceholder{Figure 55: PXL Structure}{PXL 结构}{fig:pxl-structure}

% ============================================================
\section{一般定义}
\subsection*{General Definitions}

本节给出 PXL 中部分语法元素的定义与特征，例如标识符、字符串字面量、关键字、运算符与注释。理解这些语法元素的用途与作用，有助于编写 runset。

\subsection{标识符}
\subsubsection*{Identifiers}

任意 PXL 组件的标识符均可按任意顺序包含字母字符、数字与下划线。
PXL 区分大小写。下列标识符彼此不同：
\begin{itemize}
  \item \cmd{abCD}
  \item \cmd{ABCD}
  \item \cmd{abcd}
\end{itemize}

标识符须遵守下列命名限制：
\begin{itemize}
  \item 双下划线前缀（\cmd{\_\_}）。不允许以两个下划线开头的标识符；保留给语言内部使用。
  \item 单下划线前缀（\cmd{\_}）。不允许以一个下划线开头的标识符；保留给产品使用。
  \item 全大写。保留给应用程序编程接口（API）使用。不过，可将全大写标识符用作枚举类型元素。
  \item 保留关键字。不得将保留关键字用作标识符。PXL 保留关键字列表见第~\ref{sec:pxl-keywords}~节「关键字」。
\end{itemize}

下列名称也不可用作标识符：
\begin{itemize}
  \item 函数名：这些函数定义于 \textit{IC Validator Reference Manual} 与 IC Validator 头文件中。示例函数名包括 \cmd{capacitor}、\cmd{check\_property}、\cmd{internal1} 与 \cmd{xor}。
  \item 容器类型名：列表、哈希与结构体的名称定义于 IC Validator 头文件中。容器类型名示例：
  \begin{itemize}
    \item 列表：\cmd{coordinate\_l} 用于定义 \cmd{property()} 函数
    \item 哈希：\cmd{layer\_list\_h} 用于定义 \cmd{density()} 函数
    \item 结构体：\cmd{check\_property\_s} 用于定义 \cmd{check\_property()} 函数
  \end{itemize}
\end{itemize}

合法标识符示例：
\begin{lstlisting}[style=pxl]
This
IsAnIdentifier
another_identifier_using_underscores
layer20
\end{lstlisting}

保留（非法）标识符示例：
\begin{lstlisting}[style=pxl]
__SYSTEMNAME
\end{lstlisting}
（双下划线前缀；保留给语言内部使用）

\begin{lstlisting}[style=pxl]
_productname
\end{lstlisting}
（单下划线前缀；保留给产品使用）

\begin{lstlisting}[style=pxl]
HIERARCHICAL
\end{lstlisting}
（保留关键字）

\subsection{字符串字面量}
\subsubsection*{String Literals}

字符串字面量是以双引号（\cmd{"}）开始并结束的一组字符。

\begin{noteBox}
文本字符串的第一个字符不能是整数，也不能是后跟整数的空格。\\
字符串字面量不得包含换行。若字符串字面量中需要换行，必须使用换行字符（\cmd{\textbackslash n}）。
\end{noteBox}

下列字符组均为字符串字面量示例：
\begin{lstlisting}[style=pxl]
"hello"
"\""
"a string with\nmultiple lines"
"abc" "def"
\end{lstlisting}

字符串字面量会自动与相邻的字符串字面量合并（字符串拼接）。尤其是，它们会与包含字符串值的预处理器变量（如 \cmd{\_\_FILE\_\_}）以及环境变量引用（如 \cmd{\$<id>}）合并。

可在字符串内使用转义序列，将部分字符串内容标记为给编译器的指令。PXL 可用的转义序列见表~\ref{tab:pxl-escape}。

\begin{longtable}{@{}>{\raggedright\arraybackslash\ttfamily}p{0.22\textwidth} p{0.70\textwidth}@{}}
\caption{转义序列（Escape Sequences）}\label{tab:pxl-escape}\\
\toprule
\textbf{\textrm{转义}} & \textbf{\textrm{定义}} \\
\midrule
\endfirsthead
\toprule
\textbf{\textrm{转义}} & \textbf{\textrm{定义}} \\
\midrule
\endhead
\bottomrule
\endfoot
\textbackslash n & 换行字符 \\
\textbackslash" & 双引号字符 \\
\textbackslash\textbackslash & 单个反斜杠字符 \\
\textbackslash NNN & 字符的八进制值 \\
\textbackslash t & 制表符 \\
\textbackslash xNN & 字符的十六进制值 \\
\end{longtable}

\subsection{数值常量}
\subsubsection*{Numeric Constants}

PXL 中的数值常量定义见表~\ref{tab:pxl-numeric}。

\begin{longtable}{@{}p{0.26\textwidth} p{0.42\textwidth} >{\raggedright\arraybackslash\ttfamily}p{0.24\textwidth}@{}}
\caption{数值常量（Numeric Constants）}\label{tab:pxl-numeric}\\
\toprule
\textbf{数值类型} & \textbf{条件} & \textbf{\textrm{示例}} \\
\midrule
\endfirsthead
\toprule
\textbf{数值类型} & \textbf{条件} & \textbf{\textrm{示例}} \\
\midrule
\endhead
\bottomrule
\endfoot
double & 含指数 & 4e5, 10.305e+10 \\
double & 含小数点 & 3.141592 \\
十进制整数（decimal integer） & 以数字 1--9 开头 & 1023 \\
八进制数（octal number） & 以零（0）开头，且各位为 0--7 & 0346 \\
十六进制数（hexadecimal number） & 以 \cmd{0x} 或 \cmd{0X} 开头，且仅含十六进制数字（0--9、a--f、A--F） & 0xDEADBEEF \\
\end{longtable}

\subsection{关键字}
\subsubsection*{Keywords}
\label{sec:pxl-keywords}

关键字是语言识别的内置词。关键字为保留字，因此不能用作标识符。

关键字的若干用途包括：
\begin{itemize}
  \item 流程控制，如 \cmd{if}、\cmd{foreach} 与 \cmd{while}
  \item 运算符，如 \cmd{in} 与 \cmd{contains}
  \item 常量，如 \cmd{true} 与 \cmd{false}
\end{itemize}

下列各表列出 PXL 中的保留字。

\begin{noteBox}
有关命名限制的更多信息，见「标识符」小节。
\end{noteBox}

表~\ref{tab:pxl-sys-kw} 列出由系统保留、外部不可用的关键字。

\begin{longtable}{@{}>{\raggedright\arraybackslash\ttfamily}p{0.30\textwidth} >{\raggedright\arraybackslash\ttfamily}p{0.30\textwidth} >{\raggedright\arraybackslash\ttfamily}p{0.30\textwidth}@{}}
\caption{系统关键字（System Keywords）}\label{tab:pxl-sys-kw}\\
\toprule
\endfirsthead
\endhead
\bottomrule
\endfoot
@@SYS_KW@@
\end{longtable}

表~\ref{tab:pxl-indesign-kw} 列出 In-Design 流程使用的关键字。请勿在 IC Validator runset 中使用这些名称作为变量。

\begin{longtable}{@{}>{\raggedright\arraybackslash\ttfamily}p{0.48\textwidth} >{\raggedright\arraybackslash\ttfamily}p{0.48\textwidth}@{}}
\caption{In-Design 流程关键字（In-Design Flow Keywords）}\label{tab:pxl-indesign-kw}\\
\toprule
\endfirsthead
\endhead
\bottomrule
\endfoot
@@INDESIGN@@
\end{longtable}

表~\ref{tab:pxl-app-ids} 列出不可用作标识符的保留应用程序标识符。

\begin{longtable}{@{}>{\raggedright\arraybackslash\ttfamily}p{0.30\textwidth} >{\raggedright\arraybackslash\ttfamily}p{0.30\textwidth} >{\raggedright\arraybackslash\ttfamily}p{0.30\textwidth}@{}}
\caption{保留应用程序标识符（Reserved Application Identifiers）}\label{tab:pxl-app-ids}\\
\toprule
\endfirsthead
\endhead
\bottomrule
\endfoot
@@APP_IDS@@
\end{longtable}

枚举类型元素也不能用作标识符。表~\ref{tab:pxl-enum-elems} 列出 IC Validator 工具使用的枚举类型元素。

\begin{longtable}{@{}>{\raggedright\arraybackslash\ttfamily}p{0.30\textwidth} >{\raggedright\arraybackslash\ttfamily}p{0.30\textwidth} >{\raggedright\arraybackslash\ttfamily}p{0.30\textwidth}@{}}
\caption{枚举类型元素（Enumerated Type Elements）}\label{tab:pxl-enum-elems}\\
\toprule
\endfirsthead
\endhead
\bottomrule
\endfoot
@@ENUM@@
\end{longtable}

\subsection{运算符}
\subsubsection*{Operators}

运算符是具有特定含义的符号。例如，\cmd{+} 表示将两个数相加。运算符类型见表~\ref{tab:pxl-op-types}。

\begin{longtable}{@{}p{0.22\textwidth} p{0.28\textwidth} p{0.42\textwidth}@{}}
\caption{运算符类型（Operator Types）}\label{tab:pxl-op-types}\\
\toprule
\textbf{运算符类型} & \textbf{操作数个数} & \textbf{示例} \\
\midrule
\endfirsthead
\toprule
\textbf{运算符类型} & \textbf{操作数个数} & \textbf{示例} \\
\midrule
\endhead
\bottomrule
\endfoot
一元（unary） & 1 & 逻辑非（\cmd{!}） \\
二元（binary） & 2 & 加法（\cmd{+}）、幂运算（\cmd{**}） \\
三元（ternary） & 3 & 条件运算符（\cmd{e1 ? e2 : e3}） \\
\end{longtable}

\subsection{注释}
\subsubsection*{Comments}

注释是用于说明程序的文本，便于向读者表明代码意图。注释按空白处理。表~\ref{tab:pxl-comment-types} 按注释所占行数给出注释类型。

\begin{longtable}{@{}p{0.22\textwidth} >{\raggedright\arraybackslash\ttfamily}p{0.18\textwidth} >{\raggedright\arraybackslash\ttfamily}p{0.22\textwidth} p{0.28\textwidth}@{}}
\caption{注释类型（Comment Types）}\label{tab:pxl-comment-types}\\
\toprule
\textbf{\textrm{注释类型}} & \textbf{\textrm{起始}} & \textbf{\textrm{结束}} & \textbf{\textrm{长度}} \\
\midrule
\endfirsthead
\toprule
\textbf{\textrm{注释类型}} & \textbf{\textrm{起始}} & \textbf{\textrm{结束}} & \textbf{\textrm{长度}} \\
\midrule
\endhead
\bottomrule
\endfoot
\textrm{行注释（line）} & // & 下一换行 & 一行 \\
\textrm{块注释（block）} & /* & */ & 任意长度 \\
\end{longtable}

块注释不可嵌套。若在注释内看到 \cmd{*/}，预处理器会报告警告。例如，若 runset 含有：
\begin{lstlisting}[style=pxl]
/*           - 1st comment start
       /*       - 2nd comment start
       */       - 2nd comment end
*/           - 1st comment end
\end{lstlisting}

则「1st comment start」与「2nd comment end」配对；「1st comment start」不会与「1st comment end」配对。

注释示例如下：
\begin{lstlisting}[style=pxl]
/* this is a single-line block comment */

// this is a line comment, so /* is ignored

/* block comments /* cannot nest. Therefore this comment evokes a
warning from the preprocessor */

/*
 * block comments can span
 * several lines, such
 * as this one does.
 */

/*
  A leading * is not required on each line of a
  block comment but does make the comment stand out
  from the program.
*/
\end{lstlisting}

\subsection{环境变量引用}
\subsubsection*{Environment Variable Reference}

环境变量是一组可用于控制进程运行方式的动态值。

环境变量引用由美元符（\cmd{\$}）后跟任意合法标识符组成。若引用未定义的环境变量，将在编译期报错并拒绝。

环境变量会被替换为包含其值的字符串字面量。仅对形如 \cmd{\$<id>} 的记号进行替换。具体而言，字符串字面量内部不会替换；由于相邻字符串字面量会拼接，通常也无需在字面量内替换。

\begin{noteBox}
若环境变量的值为空字符串，则将其转换为值为空字符串的字符串字面量。
\end{noteBox}

环境变量用法示例如下：
\begin{lstlisting}[style=pxl]
x : string = $y;

/*
 * assuming environment variable my_value contains the value abcd,
 * the following code evaluates to true; otherwise evaluates to false.
 */
res : boolean = ( $my_value == "abcd" );

/*
 * if environment variable aaa contains the value
 * /home/dir/path, the following code produces a string
 * containing a path to the file "cases" in that directory.
 */
path : string = $aaa "/cases";

/*
 * if environment variable bbb is defined but empty, the
 * following code produces the string abcdef.
 */
mystr : string = "abc" $bbb "def";
\end{lstlisting}
